% ═══════════════════════════════════════════════════════════════
% MT-06c-ML: Minitest MUSTERLÖSUNG - Las direcciones
% Spanisch Klasse 7 | Sequenz 6: Mi ciudad
% ═══════════════════════════════════════════════════════════════

\documentclass[11pt, a4paper]{article}

\usepackage[utf8]{inputenc}
\usepackage[T1]{fontenc}
\usepackage[ngerman]{babel}
\usepackage{helvet}
\renewcommand{\familydefault}{\sfdefault}

\usepackage[margin=2cm]{geometry}
\usepackage{enumitem}
\usepackage{tabularx}
\usepackage{booktabs}

\usepackage{xcolor}
\usepackage{amssymb}
\usepackage{tcolorbox}
\tcbuselibrary{skins}

\usepackage{fancyhdr}
\usepackage{lastpage}
\usepackage{needspace}

% FARBEN
\definecolor{spanisch-color}{HTML}{c41e3a}
\definecolor{grau-color}{HTML}{888888}
\definecolor{hellgrau-color}{HTML}{f5f5f5}
\definecolor{loesung-color}{HTML}{c62828}

\colorlet{spanisch}{spanisch-color}
\colorlet{grau}{grau-color}
\colorlet{hellgrau}{hellgrau-color}
\colorlet{loesung}{loesung-color}

\newcommand{\fachfarbe}{spanisch}

\newcommand{\thema}{Minitest: Las direcciones}
\newcommand{\fach}{Spanisch}
\newcommand{\klasse}{7}
\newcommand{\maxpunkte}{26}
\newcommand{\zeit}{15 Minuten}
\newcommand{\datum}{\today}

\pagestyle{fancy}
\fancyhf{}
\renewcommand{\headrulewidth}{0.4pt}
\renewcommand{\footrulewidth}{0.4pt}

\lhead{\fach{} -- Klasse \klasse}
\chead{\textcolor{loesung}{\textbf{MT-06c MUSTERLÖSUNG}}}
\rhead{\datum}

\lfoot{\zeit}
\cfoot{}
\rfoot{Seite \thepage\ von \pageref{LastPage}}

\newcommand{\punktebox}[1]{\hfill\fbox{\small \_\_\_/#1}}
\newcommand{\luecke}[1]{\rule{#1}{0.4pt}}
\newcommand{\lsg}[1]{\textcolor{loesung}{\textbf{#1}}}

\newtcolorbox{infobox}{
    colback=hellgrau,
    colframe=\fachfarbe,
    boxrule=1pt,
    arc=0pt,
    left=8pt, right=8pt, top=8pt, bottom=8pt
}

\newenvironment{aufgabe}[1]{%
    \needspace{5\baselineskip}%
    \nopagebreak[4]%
    \vspace{10pt}
    \noindent\textbf{\textcolor{\fachfarbe}{Aufgabe #1}}
    \nopagebreak[4]%
    \vspace{5pt}
    \nopagebreak[4]%
    \begin{list}{}{\leftmargin=0pt}
    \item[]
}{%
    \end{list}
}

\newcommand{\es}[1]{\textcolor{\fachfarbe}{\textbf{#1}}}

\begin{document}

% ═══════════════════════════════════════════════════════════════
% KOPFBEREICH
% ═══════════════════════════════════════════════════════════════

\begin{center}
    {\Large\bfseries\textcolor{\fachfarbe}{\thema}}\\[0.3em]
    {\large\textcolor{loesung}{\textbf{--- MUSTERLÖSUNG ---}}}
\end{center}

\vspace{5pt}

\noindent
\begin{tabularx}{\textwidth}{|X|X|X|}
\hline
\textbf{Name:} \lsg{Musterschüler/in} &
\textbf{Klasse:} \klasse &
\textbf{Datum:} \datum \\
\hline
\end{tabularx}

\vspace{10pt}

\begin{infobox}
\textbf{Zeit:} \zeit{} | \textbf{Punkte:} \maxpunkte{} BE | \textbf{Hilfsmittel:} keine
\end{infobox}

\vspace{15pt}

% ═══════════════════════════════════════════════════════════════
% AUFGABE 1: Richtungsangaben zuordnen
% ═══════════════════════════════════════════════════════════════

\begin{aufgabe}{1 -- Richtungsangaben zuordnen}
Verbinde die spanischen Ausdrücke mit der deutschen Bedeutung. \punktebox{6}
\end{aufgabe}

\begin{center}
\begin{tabular}{rcl}
$\uparrow$ \es{todo recto} & \lsg{---} & nach rechts \lsg{$\rightarrow$ geradeaus} \\[0.8em]
$\rightarrow$ \es{a la derecha} & \lsg{---} & nach links \lsg{$\rightarrow$ nach rechts} \\[0.8em]
$\leftarrow$ \es{a la izquierda} & \lsg{---} & geradeaus \lsg{$\rightarrow$ nach links} \\[0.8em]
\es{girar} & \lsg{---} & abbiegen \lsg{$\checkmark$} \\[0.8em]
\es{seguir} & \lsg{---} & weitergehen \lsg{$\checkmark$} \\[0.8em]
\es{cruzar} & \lsg{---} & überqueren \lsg{$\checkmark$} \\
\end{tabular}
\end{center}

\vspace{15pt}

% ═══════════════════════════════════════════════════════════════
% AUFGABE 2: Ortsangaben zuordnen
% ═══════════════════════════════════════════════════════════════

\begin{aufgabe}{2 -- Ortsangaben zuordnen}
Verbinde die spanischen Ortsangaben mit der deutschen Bedeutung. \punktebox{6}
\end{aufgabe}

\begin{center}
\begin{tabular}{rcl}
\es{al lado de} & \lsg{---} & gegenüber von \lsg{$\rightarrow$ neben} \\[0.8em]
\es{enfrente de} & \lsg{---} & neben \lsg{$\rightarrow$ gegenüber von} \\[0.8em]
\es{cerca de} & \lsg{---} & weit weg von \lsg{$\rightarrow$ in der Nähe von} \\[0.8em]
\es{lejos de} & \lsg{---} & in der Nähe von \lsg{$\rightarrow$ weit weg von} \\[0.8em]
\es{entre ... y ...} & \lsg{---} & hinter \lsg{$\rightarrow$ zwischen ... und ...} \\[0.8em]
\es{detrás de} & \lsg{---} & zwischen ... und ... \lsg{$\rightarrow$ hinter} \\
\end{tabular}
\end{center}

\vspace{15pt}

% ═══════════════════════════════════════════════════════════════
% AUFGABE 3: Lückentext
% ═══════════════════════════════════════════════════════════════

\begin{aufgabe}{3 -- Sätze ergänzen}
Ergänze die Lücken mit dem passenden spanischen Ausdruck. \punktebox{8}
\end{aufgabe}

\noindent
a) La farmacia está \lsg{enfrente de} la panadería. (gegenüber)

\vspace{0.8em}

\noindent
b) El banco está \lsg{al lado de/del} el supermercado. (neben)

\vspace{0.8em}

\noindent
c) Sigue \lsg{todo recto} y llegas a la plaza. (geradeaus)

\vspace{0.8em}

\noindent
d) La estación está \lsg{lejos} de aquí. (weit weg)

\vspace{0.8em}

\noindent
e) El museo está \lsg{cerca} del parque. (in der Nähe)

\vspace{0.8em}

\noindent
f) Gira \lsg{a la derecha}. (nach rechts)

\vspace{0.8em}

\noindent
g) El cine está \lsg{detrás de} la biblioteca. (hinter)

\vspace{0.8em}

\noindent
h) Gira \lsg{a la izquierda}. (nach links)

\vspace{15pt}

% ═══════════════════════════════════════════════════════════════
% AUFGABE 4: Fragen stellen
% ═══════════════════════════════════════════════════════════════

\begin{aufgabe}{4 -- Fragen stellen}
Wie fragst du auf Spanisch? Kreise die richtige Frage ein. \punktebox{6}
\end{aufgabe}

\noindent
a) ``Wo ist die Apotheke?''

\hspace{1cm} A) ¿Cómo llego a la farmacia? \quad \lsg{B) ¿Dónde está la farmacia?} \quad C) ¿Qué es la farmacia?

\vspace{0.8em}

\noindent
b) ``Wie komme ich zum Bahnhof?''

\hspace{1cm} A) ¿Dónde está la estación? \quad B) ¿Qué es la estación? \quad \lsg{C) ¿Cómo llego a la estación?}

\vspace{0.8em}

\noindent
c) ``Ist es weit?''

\hspace{1cm} \lsg{A) ¿Está lejos?} \quad B) ¿Está cerca? \quad C) ¿Dónde está?

\vspace{0.8em}

\noindent
d) ``Ist es in der Nähe?''

\hspace{1cm} A) ¿Está lejos? \quad \lsg{B) ¿Está cerca?} \quad C) ¿Cómo llego?

\vspace{0.8em}

\noindent
e) ``Wo ist das Museum?''

\hspace{1cm} A) ¿Cómo llego al museo? \quad \lsg{B) ¿Dónde está el museo?} \quad C) ¿Qué es el museo?

\vspace{0.8em}

\noindent
f) ``Wie komme ich zur Bibliothek?''

\hspace{1cm} A) ¿Dónde está la biblioteca? \quad \lsg{B) ¿Cómo llego a la biblioteca?} \quad C) ¿Qué es?

% ═══════════════════════════════════════════════════════════════
% PUNKTEÜBERSICHT
% ═══════════════════════════════════════════════════════════════

\vspace{25pt}
\hrule
\vspace{10pt}

\noindent
\begin{tabularx}{\textwidth}{|X|c|c|c|c||c|}
\hline
\textbf{Aufgabe} & 1 & 2 & 3 & 4 & \textbf{Gesamt} \\
\hline
\textbf{max. Punkte} & 6 & 6 & 8 & 6 & \textbf{26} \\
\hline
\textbf{erreicht} & \lsg{6} & \lsg{6} & \lsg{8} & \lsg{6} & \lsg{26} \\
\hline
\end{tabularx}

\vspace{10pt}

\begin{flushright}
\textbf{Note:} \lsg{14 NP}
\end{flushright}

\vspace{1em}

\textbf{Notenspiegel (14-NP-Skala):}

\begin{center}
\footnotesize
\begin{tabular}{|c|c|c|c|c|c|c|c|c|c|c|c|c|c|c|}
\hline
\textbf{NP} & 14 & 13 & 12 & 11 & 10 & 9 & 8 & 7 & 6 & 5 & 4 & 3 & 2 & 1 \\
\hline
\textbf{ab BE} & 26 & 25 & 24 & 22 & 21 & 19 & 17 & 16 & 14 & 12 & 10 & 9 & 7 & 5 \\
\hline
\end{tabular}
\end{center}

\vspace{1em}

\textcolor{loesung}{\textbf{Merke:}}

\textcolor{loesung}{\textit{¿Dónde está...?} = Wo ist...? (Lage)}

\textcolor{loesung}{\textit{¿Cómo llego a...?} = Wie komme ich zu...? (Weg)}

\end{document}
